% Supporting-information outline for implementation details moved out of the main text.
% This file is not yet assembled into the SI PDF. It records the intended SI content
% so that the main manuscript can remain physics-first.

\section{Implementation details for the detector-space calculation}

The main text describes the calculation in reciprocal-space and detector-space language. The numerical implementation uses $2\theta$--$\phi$ remapping and caked intensity fields for efficiency, but those coordinates do not change the physical comparison. The same reciprocal-space trajectories and projection regions are applied to the measured and calculated images.

\subsection{Mosaic surface density}

The out-of-plane mosaic line density $\omega(\Delta\theta)$ is evaluated as a wrapped Gaussian-plus-Lorentzian distribution. After random in-plane rotation about the film normal, the line density becomes an axially symmetric density on the Bragg sphere. For a reflection with $G=|\mathbf G|$, polar angle $\vartheta$, and Bragg position $\vartheta_B$, the normalized surface density used for ray weights is
\begin{equation}
\rho(\vartheta)=\frac{I_0}{2\pi G^2\sin\vartheta}\,\omega(\vartheta-\vartheta_B),\qquad 0<\vartheta<\pi.
\end{equation}
The $1/\sin\vartheta$ term is the Jacobian for converting the one-dimensional tilt density into a surface density.

\subsection{Coordinate remapping and binning}

Detector pixels are transformed into $2\theta$--$\phi$ fields with sub-pixelation so that curved fixed-$m$ trajectories and projected $Q_z$ profiles can be integrated without pixel-boundary artifacts. Lookup tables connect the detector grid and the remapped fields during fitting.

\subsection{Monte Carlo sampling}

Beam position, divergence, wavelength spread, and mosaic events are sampled probabilistically for computational efficiency. This replaces a dense grid over the same phase space without changing the physical model used to generate the detector image.


\subsection{Projection and intensity uncertainty}

The transformed-coordinate uncertainty is evaluated from the detector-pixel footprint rather than assigned after the caked or projected image is produced. For square detector pixels with pitch $p$ and detector distance $d$, the leading flat-detector terms are
\begin{equation}
\sigma_{2\theta,\mathrm{pix}}(t)=\frac{p\cos^{2}t}{\sqrt{12}\,d},
\qquad
\sigma_{\phi,\mathrm{pix}}(t)=\frac{p}{\sqrt{12}\,d\tan t},
\qquad t=2\theta .
\end{equation}
The radial term is finite and largest near the direct beam. The azimuthal term scales as $1/\tan(2\theta)$, so the problematic regime is confined to the near-specular, low-$2\theta$ region where azimuth is poorly conditioned.

For the $d=75$ mm and $p=300\,\mu\mathrm{m}$ geometry used in the uncertainty calculation, the maximum radial coordinate uncertainty is approximately $0.066^{\circ}$. The azimuthal coordinate uncertainty falls below $1^{\circ}$ above $2\theta\approx3.8^{\circ}$, below $0.5^{\circ}$ above $2\theta\approx7.5^{\circ}$, and below $0.25^{\circ}$ above $2\theta\approx14.9^{\circ}$. More generally, the threshold for a target azimuthal width $w_{\phi}$ is
\begin{equation}
2\theta_{\mathrm{crit}}=\arctan\!\left(\frac{p}{\sqrt{12}\,d\,w_{\phi}}\right),
\end{equation}
with $w_{\phi}$ expressed in radians.

The intensity uncertainty of a projected bin is propagated from detector-space variance through the same overlap weights used to construct the rebinned image. If $s_i$, $n_i$, and $v_i$ are the signal, normalization, and variance of detector pixel $i$, and $w_{ib}$ is the overlap weight from detector pixel $i$ to projected bin $b$, then
\begin{equation}
S_b=\sum_i w_{ib}s_i,
\qquad
N_b=\sum_i w_{ib}n_i,
\qquad
I_b=\frac{S_b}{N_b},
\end{equation}
with
\begin{equation}
\mathrm{Var}(I_b)=\frac{\sum_i w_{ib}^{2}v_i}{N_b^{2}}.
\end{equation}
Thus the profile uncertainty remains detector-grounded after the coordinate transformation. This uncertainty propagation supports the interpretation of the projected line shapes, but it is not a complete posterior covariance analysis for every fitted model parameter.
